%!TEX root = ../thesis.tex
\chapter*{Abstract}
\pagestyle{pre}

Wind turbines are equipped with sensors that record many different measurements every
ten minutes. Making sense of these measurements is not straightforward, because a wind
turbine behaves very differently depending on what it is doing at the time of the
recording, whether it is in standby, ramping up, producing power at partial load, or
running at full capacity. Without knowing which of these operating states the turbine
is currently in, it is difficult to build reliable models for fault detection, power
curve analysis, or maintenance planning.

This thesis develops and evaluates methods for classifying wind turbine operating
states, referred to as operational zones, from SCADA sensor data. The work is structured
as two connected studies. The first study defines nine operational zones based on turbine
physics and trains rule-based classifiers (RIPPER and FOIL) to identify these zones
automatically. The second study takes the same zones and compares deep reinforcement
learning (DRL) agents (DQN, PPO, A2C) with supervised machine learning classifiers
(Random Forest, Gradient Boosting, MLP) on the same task.

A key finding of the second study is that a Random Forest trained on only four raw SCADA
signals achieves 100\% accuracy on both the 2016 validation set and an independent 2017
test set. This result confirms that the zone definitions are not only physically
meaningful but also cleanly separable in the raw sensor space. Among the DRL agents,
DQN with physics guardrails reaches 97.7\% accuracy on the external test, a strong
result for a learning-based approach that does not require labelled training data in
principle. A deployment pipeline comparison shows that placing DQN before a LightGBM
meta-model reduces the number of readings flagged for human review by approximately
70\% compared to building a meta-model directly on RIPPER. Together, the two studies
provide practical guidance on which approach is most suitable depending on whether
labelled data is available, interpretability is required, or online adaptation is needed.

\bigskip
\noindent\textbf{Keywords:} wind turbine, SCADA, operational zone classification,
prognostics and health management, digital twin, rule-based machine learning, RIPPER,
deep reinforcement learning, Random Forest
